% Mock-ups of the MLST paper's main-text figures and tables, in the layouts of the figure plan
% (docs/FIGURE_PLAN.md, local), on the data committed today. Not part of the paper: the figures
% are drawn by experiments/FIGS/mockup_figures.py, and every one carries "[preliminary data]".
% Compiled by the render job (scripts/build_render_job.py --mockups).
\documentclass[11pt]{article}
\usepackage[a4paper,textwidth=15cm,top=2.2cm,bottom=2.2cm]{geometry}
\usepackage{amsmath,amssymb}
\usepackage{booktabs}
\usepackage{graphicx}
\usepackage{xcolor}
\usepackage{placeins}
\usepackage[font=small,labelfont=bf]{caption}
\graphicspath{{../../figures/mockups/}}
\newcommand{\prelim}{\textcolor{red!70!black}{[preliminary data]}}
\newcommand{\e}{$\cdot$}
\setlength{\parindent}{0pt}
\setlength{\parskip}{4pt}

\title{Label granularity in supervised pretraining of jet foundation models\\[4pt]
\large Main-text figure and table mock-ups \prelim}
\date{2 November 2026}
\begin{document}
\maketitle

Each item is drawn in the layout the figure plan gives it, from the results committed today.
Each caption gives the claim the item will carry once the final results are in. The note under each figure
says what is a stand-in or old. Tables are skeletons: rows and columns as planned, with
cells left empty (\e) until the final numbers are generated, except Table~1, which is final.
Paper order: Figure~1, Table~1, Figures~2--4, Table~2, Figure~5, Tables~3--4, Figure~6 (F7),
Figure~7 (F8a), Table~5, Figure~8 (F8b), Figure~9. The benchmark ROC figure (F6) is left out
(decision P3).

% ------------------------------------------------------------------ F1
\begin{figure}[p]
\centering
\begin{minipage}[t]{0.55\linewidth}\vspace{0pt}
\includegraphics[width=\linewidth]{label_sets.pdf}\\
\small (a) Class hierarchy, with a lineage of decays through the four levels and the controls.
\end{minipage}\hfill
\begin{minipage}[t]{0.42\linewidth}\vspace{0pt}
\fbox{\parbox{0.95\linewidth}{\small (b) \textbf{Controls} (TikZ in the final figure).
A random partition: resonant classes assigned at random to 16 groups plus QCD, group sizes
matched to the 17-class groups. Three partitions that differ only in a flavor boundary: none,
a $b$-quark boundary in four-quark decays, a random boundary of the same size.}}\\[6pt]
\fbox{\parbox{0.95\linewidth}{\small (c) \textbf{Pipeline}. Same jets, order and initial weights
for every model of a seed $\to$ pretraining on one label set (optionally with the auxiliary
mass regression) $\to$ frozen features $\to$ linear probe (MLP probe beside a null) or
fine-tuning $\to$ JetClass-II, JetClass, top, quark/gluon (Pythia, Herwig), $X\to b\bar b$/$X\to bs$.}}
\end{minipage}
\caption{\prelim\ Design. Within a seed, models see the same jets in the same order from the same
initial weights; only the label set changes, and each merge removes a named physics distinction.
The random and flavor-controlled partitions separate the number of classes from which classes
are merged. Panel (a) is today's schematic; (b) and (c) are placeholders.}
\end{figure}

% ------------------------------------------------------------------ T1
\begin{table}[p]
\centering
\caption{Analysis discriminants built from sums of Sophon's output scores, and whether each label
set can form them ($\checkmark$) or not (---). Final.}
\begin{tabular}{llcccc}
\toprule
discriminant & source & 188 & 162 & 43 & 17 \\
\midrule
$X\to b\bar b$ vs QCD & Sophon, Eq.~(4) & $\checkmark$ & $\checkmark$ & --- & --- \\
$X\to c\bar s$ vs QCD (Sophon's $A_1$) & Sophon, Eqs.~(5)--(6) & $\checkmark$ & $\checkmark$ & --- & --- \\
three quarks, exactly one $b$ (Sophon's $A_3$) & Sophon, Eqs.~(5)--(6) & $\checkmark$ & $\checkmark$ & --- & --- \\
$W'\to W\phi\to WWW$ & Sophon, Eqs.~(7)--(9) & $\checkmark$ & $\checkmark$ & --- & --- \\
$D_{bc}$, boosted $|V_{cb}|$ & Zhao et al., Eq.~(1) & $\checkmark$ & $\checkmark$ & --- & --- \\
\bottomrule
\end{tabular}
\end{table}

% ------------------------------------------------------------------ F2
\begin{figure}[p]
\centering
\includegraphics[width=0.7\linewidth]{F2_label_granularity.pdf}
\caption{\prelim\ Linear probe on the 188 JetClass-II classes: ratio of $1-\mathrm{AUC}$ (one class
against the rest, averaged) to the 188-class model of the same seed, against the number of
pretraining classes. Claim: coarser labels give features from which a linear probe separates
the classes worse, and random partitions at 17 classes lose less than the 17-class level of the
hierarchy, so the cost depends on which classes are merged. The final caption gives the
absolute AUC of the 188-class model and of the randomly initialized network, and states that this
is the 188-class model's own pretraining task. Form of Hong et al.\ Fig.~5.}
\end{figure}

% ------------------------------------------------------------------ F3
\begin{figure}[p]
\centering
\includegraphics[height=0.86\textheight]{F3_per_class.pdf}
\caption{\prelim\ JetClass-II, class by class: ratio of a class's $1-\mathrm{AUC}$ under each label set
to the 188-class model of the same seed (linear probe, mean over runs). Rows are grouped by the
17-class parent; QCD forms the bottom block. Claim: the cost concentrates in classes whose label merges
them with decays that differ only in flavor, and the controls move it to whichever classes
they merge. New form; nearest precedents JP-JEPA Fig.~2f,g and CMS JME-25-001 Fig.~7.}
\end{figure}

% ------------------------------------------------------------------ F4
\begin{figure}[p]
\centering
\includegraphics[width=\linewidth]{F4_flavor_pairs.pdf}
\caption{\prelim\ Decays that differ only in flavor: ratio of $1-\mathrm{AUC}$ on the pair to the
188-class model of the same seed. Question answered: does a flavor boundary in the label set
improve the linear separation of such pairs, including pairs the label set merges? Open
markers in the right panel: MLP probe beside the 162-class linear null. The sentence is written
from the final numbers; today's data do not support ``boundary, not count'' as phrased.}
\end{figure}

% ------------------------------------------------------------------ T2
\begin{table}[p]
\centering
\small
\caption{\prelim\ JetClass in the form of ParT's Table~1 (skeleton). Accuracy, AUC averaged over class pairs,
and the rejection of QCD jets per signal class at ParT's working points. Literature rows at
$2\times10^6$ training jets, from scratch; ours fine-tuned on $10^6$. Mean $\pm$ SD over three
pretraining runs; $^\dagger$: fewer than 20 QCD jets pass. Claim: the 162-class model matches
the 188-class model; at 17 classes rejection falls in most signal classes while $H\to b\bar b$
is unchanged.}
\setlength{\tabcolsep}{3pt}
\resizebox{\linewidth}{!}{%
\begin{tabular}{lccccccccccc}
\toprule
 & Acc & AUC & $H\to b\bar b$ & $H\to c\bar c$ & $H\to gg$ & $H\to 4q$ & $H\to\ell\nu qq'$ & $t\to bqq'$ & $t\to b\ell\nu$ & $W\to qq'$ & $Z\to q\bar q$ \\
 & & & $R_{50}$ & $R_{50}$ & $R_{50}$ & $R_{50}$ & $R_{99}$ & $R_{50}$ & $R_{99.5}$ & $R_{50}$ & $R_{50}$ \\
\midrule
ParticleNet & \e & \e & \e & \e & \e & \e & \e & \e & \e & \e & \e \\
ParT & \e & \e & \e & \e & \e & \e & \e & \e & \e & \e & \e \\
MIParT & \e & \e & \e & \e & \e & \e & \e & \e & \e & \e & \e \\
L-GATr & \e & \e & \e & \e & \e & \e & \e & \e & \e & \e & \e \\
\midrule
from scratch & \e & \e & \e & \e & \e & \e & \e & \e & \e & \e & \e \\
188 classes, fine-tuned & \e & \e & \e & \e & \e & \e & \e & \e & \e & \e & \e \\
162 classes, fine-tuned & \e & \e & \e & \e & \e & \e & \e & \e & \e & \e & \e \\
43 classes, fine-tuned & \e & \e & \e & \e & \e & \e & \e & \e & \e & \e & \e \\
17 classes, fine-tuned & \e & \e & \e & \e & \e & \e & \e & \e & \e & \e & \e \\
\bottomrule
\end{tabular}}
\end{table}

% ------------------------------------------------------------------ F5
\begin{figure}[p]
\centering
\includegraphics[width=\linewidth]{F5_missing_inputs.pdf}
\caption{\prelim\ Missing input features. Top: relative drop in JetClass-II accuracy when the
top-tagging or quark/gluon inputs are zero (and, when available, set to the training mean or to an
average of random draws of real particles). Bottom: $R_{30}$ of the linear probe on top and
quark/gluon, open circles the MLP probe. Claim: the weak frozen probes on top and quark/gluon
come from zeroed inputs placing jets where no pretraining jet lies, not from lost information.
Form of JP-JEPA Fig.~3a with Elsharkawy et al.\ Fig.~4 hatching.}
\end{figure}

% ------------------------------------------------------------------ T3
\begin{table}[p]
\centering
\small
\caption{\prelim\ Top tagging (skeleton). Literature rows as published, cited by number; this work:
fine-tuned with 10 of 17 inputs zero, mean $\pm$ SD over three pretraining runs. Best value per
column in bold, with ours within one SD. Footnote: ParT's script on our cluster, then with
Sophon's output head, then with Sophon's jet-$p_{\mathrm T}$ rescaling (appendix). Claim: the
pretrained models improve on the same network trained from scratch and differ from each other
by about the seed spread.}
\begin{tabular}{lccccc}
\toprule
 & inputs zero & Acc & AUC & $R_{50}$ & $R_{30}$ \\
\midrule
ParticleNet & --- & \e & \e & \e & \e \\
ParT & --- & \e & \e & \e & \e \\
ParT-f.t. & --- & \e & \e & \e & \e \\
MIParT-L-f.t. & --- & \e & \e & \e & \e \\
L-GATr-f.t. & --- & \e & \e & \e & \e \\
OmniLearned-l & --- & \e & \e & \e & \e \\
\midrule
from scratch & 10 & \e & \e & \e & \e \\
188 classes, fine-tuned & 10 & \e & \e & \e & \e \\
162 classes, fine-tuned & 10 & \e & \e & \e & \e \\
43 classes, fine-tuned & 10 & \e & \e & \e & \e \\
17 classes, fine-tuned & 10 & \e & \e & \e & \e \\
188 classes, kinematic inputs only (diagnostic) & 0 & \e & \e & \e & \e \\
\bottomrule
\end{tabular}
\end{table}

% ------------------------------------------------------------------ T4
\begin{table}[p]
\centering
\small
\caption{\prelim\ Quark/gluon discrimination with ParT's ``exp'' inputs, Pythia and Herwig test sets
(skeleton). This work: fine-tuned with the four displacement inputs zero. Claim: the pretrained
models match ParT-f.t.\ on Pythia, differ from each other by about the seed spread, and lose the
same on Herwig as the network trained from scratch.}
\resizebox{\linewidth}{!}{%
\begin{tabular}{lcccccccc}
\toprule
 & \multicolumn{4}{c}{Pythia} & \multicolumn{4}{c}{Herwig} \\
\cmidrule(lr){2-5}\cmidrule(lr){6-9}
 & Acc & AUC & $R_{50}$ & $R_{30}$ & Acc & AUC & $R_{50}$ & $R_{30}$ \\
\midrule
PFN-Ex & \e & \e & \e & \e & --- & --- & --- & --- \\
ParticleNet & \e & \e & \e & \e & --- & --- & --- & --- \\
ParT & \e & \e & \e & \e & --- & --- & --- & --- \\
ParticleNet-f.t. & \e & \e & \e & \e & --- & --- & --- & --- \\
ParT-f.t. & \e & \e & \e & \e & --- & --- & --- & --- \\
OmniLearned-m & \e & \e & \e & \e & --- & --- & --- & --- \\
\midrule
from scratch & \e & \e & \e & \e & \e & \e & \e & \e \\
188 classes, fine-tuned & \e & \e & \e & \e & \e & \e & \e & \e \\
162 classes, fine-tuned & \e & \e & \e & \e & \e & \e & \e & \e \\
43 classes, fine-tuned & \e & \e & \e & \e & \e & \e & \e & \e \\
17 classes, fine-tuned & \e & \e & \e & \e & \e & \e & \e & \e \\
\bottomrule
\end{tabular}}
\end{table}

% ------------------------------------------------------------------ F7
\begin{figure}[p]
\centering
\includegraphics[width=\linewidth]{F7_finetune_vs_size.pdf}
\caption{\prelim\ Fine-tuning against the number of training jets: the 188-class model fine-tuned
(solid) and probed linearly (dashed), and the same network trained from scratch, with the ratio
to from scratch below. Claim: pretraining helps most with few training jets, and the gain
shrinks toward the full set. Form of Golling et al.\ Fig.~3, OmniJet-$\alpha$ Fig.~7 and
Vigl et al.\ Fig.~4. The other label sets are in the appendix.}
\end{figure}

% ------------------------------------------------------------------ F8a
\begin{figure}[p]
\centering
\includegraphics[width=\linewidth]{F8a_mass_regression.pdf}
\caption{\prelim\ Jet mass with the auxiliary regression target: response of all matched jets, and
resolution against jet $p_{\mathrm T}$ and generator soft-drop mass, each with the ratio to the
groomed mass. Claim: the regression narrows the response at every $p_{\mathrm T}$ and generator
mass, equally at 162 and 17 classes. Form of ATL-PHYS-PUB-2024-015 Figs.~7--8.}
\end{figure}

% ------------------------------------------------------------------ T5
\begin{table}[p]
\centering
\small
\caption{\prelim\ Jet mass from frozen features (skeleton): resolution of an MLP regression on the
frozen features. Claim: the auxiliary mass target makes the frozen features carry the jet mass,
while the label sets without it differ little. Layout of CMS-DP-2026-104 Tables~1--2.}
\begin{tabular}{lc}
\toprule
 & resolution, MLP on frozen features \\
\midrule
groomed (soft-drop) mass & \e \\
randomly initialized network & \e \\
188 classes & \e \\
162 classes & \e \\
43 classes & \e \\
17 classes & \e \\
162 classes, mass regression & \e \\
17 classes, mass regression & \e \\
\bottomrule
\end{tabular}
\end{table}

% ------------------------------------------------------------------ F8b
\begin{figure}[p]
\centering
\includegraphics[width=\linewidth]{F8b_sculpting.pdf}
\caption{\prelim\ Mass sculpting: QCD soft-drop mass before and after a cut on the summed resonance
score at 10\% and 1\% QCD efficiency, and the Jensen--Shannon distance against the QCD
efficiency. Question answered: does selecting on the tagger make the background mass peak where
a signal would? Claim: at tight selections coarser label sets reshape the QCD mass more, above
what a random selection gives. Form of CMS JME-25-001 Figs.~8--9.}
\end{figure}

% ------------------------------------------------------------------ F9
\begin{figure}[p]
\centering
\IfFileExists{../../figures/mockups/F9_xbx_standin.pdf}{%
\includegraphics[width=\linewidth]{F9_xbx_standin.pdf}}{%
\fbox{\parbox{0.8\linewidth}{\centering\vspace{2cm} Stand-in figure not found on /data at render time.\vspace{2cm}}}}
\caption{\prelim\ $X\to b\bar b$ and $X\to bs$ tagging, the form of Sophon Fig.~2: $\varepsilon_S/\sqrt{\varepsilon_B}$
against QCD efficiency. \textbf{Stand-in}: the dry run on JetClass-II jets as signal, not the
$X\to b\bar b$/$X\to bs$ samples (expected 3 November). The final figure has one panel per
signal, all four label sets, Sophon's direct discriminant on the 188- and 162-class models, the
from-scratch network on $X\to bs$, and the closest-output-node control.}
\end{figure}

\end{document}
